%%=============================================================================
%% Inleiding
%%=============================================================================

\chapter{\IfLanguageName{dutch}{Inleiding}{Introduction}}%
\label{ch:inleiding}

%De inleiding moet de lezer net genoeg informatie verschaffen om het onderwerp te begrijpen en in te zien waarom de onderzoeksvraag de moeite waard is om te onderzoeken. In de inleiding ga je literatuurverwijzingen beperken, zodat de tekst vlot leesbaar blijft. Je kan de inleiding verder onderverdelen in secties als dit de tekst verduidelijkt. Zaken die aan bod kunnen komen in de inleiding~\autocite{Pollefliet2011}:

%\begin{itemize}
%  \item context, achtergrond
%  \item afbakenen van het onderwerp
%  \item verantwoording van het onderwerp, methodologie
 % \item probleemstelling
%  \item onderzoeksdoelstelling
%  \item onderzoeksvraag
%  \item \ldots
%\end{itemize}
%\begin{figure}[H]
%\centering
%\begin{tikzpicture}[
%    node distance=2cm,
%    topic/.style={rectangle, draw, rounded corners, align=center, minimum width=3cm, minimum height=1cm},
%    main/.style={topic, fill=hogent-blue!20}
%]

%\node[main] (main) {Bachelorproef\\(z/OS System Admin)};

%\node[topic, below left=of main] (t1) {Competenties};
%\node[topic, below=of main] (t2) {Operational Tasks};
%\node[topic, below right=of main] (t3) {Frameworks};

%\draw[-] (main) -- (t1);
%\draw[-] (main) -- (t2);
%\draw[-] (main) -- (t3);

%\end{tikzpicture}
%\caption{Overzicht van topics binnen de bachelorproef}
%\end{figure}
\section{\IfLanguageName{dutch}{Probleemstelling}{Problem Statement}}%
\label{sec:probleemstelling}

Dit hoofdstuk analyseert de huidige stand van zaken binnen het domein van z/OS systeembeheer en competentieontwikkeling.
Het centrale probleem dat uit de literatuur naar voren komt, is het groeiende tekort aan gekwalificeerde mainframeprofessionals, gecombineerd met een gebrek aan gestructureerde en praktijkgerichte opleidingskaders.
Dit onderzoek richt zich daarom op het ontwikkelen van een competentieframework voor z/OS System Administrators (SysAdmin) en System Programmers (SysProg), dat zowel aansluit bij industriële noden als inzetbaar is binnen een onderwijscontext.
Dit hoofdstuk onderbouwt deze probleemstelling door bestaande technologieën, rollen en competentieraamwerken te analyseren.

De mainframe-industrie, en meer specifiek de z/OS-omgeving, vormt nog steeds een fundamentele component van de IT-architectuur binnen grote organisaties zoals banken, verzekeringsmaatschappijen en overheidsinstellingen. Hoewel cloudoplossingen en gedistribueerde systemen sterk in opmars zijn, blijft z/OS verantwoordelijk voor de verwerking van een aanzienlijk volume aan bedrijfskritische transacties. Betrouwbaarheid, beveiliging en hoge beschikbaarheid zijn daarbij kernvereisten, waardoor z/OS een cruciale rol speelt in het waarborgen van de operationele continuïteit van ondernemingen.

Binnen deze context vervullen z/OS SysAdmin en SysProg een essentiële functie in het beheer,optimalisatie en de onderhoud van de omgeving. Hun verantwoordelijkheden omvatten onder meer systeemconfiguratie, onderhoud, performantiebeheer, security-implementatie en lifecycle management. In de praktijk zijn deze taken echter vaak verspreid over meerdere gespecialiseerde teams, waarbij de exacte invulling van de rol varieert afhankelijk van de organisatie en infrastructuur. Deze variabiliteit bemoeilijkt het afbakenen van een uniform en praktijkgericht competentieprofiel dat zowel technische als niet-technische vaardigheden omvat.

Tegelijkertijd wordt de sector geconfronteerd met een structureel tekort aan gekwalificeerde mainframe-specialisten. Industrieonderzoek wijst op een vergrijzende workforce en een toenemende uitstroom van ervaren medewerkers \autocite{BMC2024}. Volgens het 2025 State of Mainframe Modernization Survey Report \autocite{Kyndryl2025} wordt deze skills gap verder versterkt door de snelle evolutie van technologieën zoals cloud, artificiële intelligentie en geavanceerde security, in combinatie met de pensionering van professionals met diepgaande mainframe-expertise. Organisaties staan hierdoor voor de uitdaging om traditionele z/OS-kennis te combineren met moderne technologiecompetenties.

Daarnaast besteden onderwijsinstellingen doorgaans beperkte aandacht aan mainframe-technologieën binnen hun IT-curricula, waardoor een kloof ontstaat tussen de verworven academische kennis en de verwachtingen vanuit de industrie. Specifiek voor instapfuncties zoals z/OS SysAdmin en SysProg rijst de vraag in welke mate bestaande competentiekaders voldoende aansluiten bij de hedendaagse praktijk en de veranderende technologische context. Het ontwikkelen of verfijnen van een praktijkgericht competentieframework kan bijdragen aan een betere afstemming tussen onderwijs, industriële verwachtingen en toekomstige workforcebehoeften, en zo het duurzaam borgen van mainframe-expertise ondersteunen.

Tijdens het onderzoek werd de oorspronkelijke focus op System Administrators (Level 1 en Level 2) uitgebreid met elementen uit de System Programmer rol, op basis van inzichten uit interviews met professionals uit de sector.

\section{\IfLanguageName{dutch}{Onderzoeksvraag}{Research question}}%
\label{sec:onderzoeksvraag}
De centrale onderzoeksvraag van deze bachelorproef luidt: “Welke competenties
moet een z/OS System Administrator en/of System Programmer bezitten om effectief te functioneren, en hoe kunnen deze competenties in een
gestructureerd framework worden georganiseerd dat direct toepasbaar is binnen
een onderwijscontext?”
Deze hoofdvraag wordt ondersteund door de volgende deelvragen:
\begin{itemize}
\item Welke technische kennis en vaardigheden zijn essentieel voor een beginnende
z/OS System Administrator/System Programmer?
\item Welke soft skills worden door de industrie als cruciaal beschouwd voor deze
rol?
\item Hoe kan een competentie framework worden vertaald naar praktische, taak
gerichte oefeningen die aansluiten bij de realiteit van het werkveld?
\item Op welke wijze kan een interactieve webapplicatie bijdragen aan een effectieve implementatie van het framework in onderwijsprogramma’s
%Wees zo concreet mogelijk bij het formuleren van je onderzoeksvraag. Een onderzoeksvraag is trouwens iets waar nog niemand op dit moment een antwoord heeft (voor zover je kan nagaan). Het opzoeken van bestaande informatie (bv. ``welke tools bestaan er voor deze toepassing?'') is dus geen onderzoeksvraag. Je kan de onderzoeksvraag verder specifiëren in deelvragen. Bv.~als je onderzoek gaat over performantiemetingen, dan 
\end{itemize}
\section{\IfLanguageName{dutch}{Onderzoeksdoelstelling}{Research objective}}%
\label{sec:onderzoeksdoelstelling}
Het doel van deze bachelorproef is het ontwikkelen van een omvattend en toekomstbestendig competentieframework voor de rol van z/OS System Administrator en System Programmers op instapniveau binnen de mainframe-industrie, specifiek gericht op een onderwijscontext. Dit framework heeft als doel een gestructureerd en praktijkgericht leerpad te bieden dat studenten voorbereidt op de operationele realiteit van een z/OS-omgeving.

Het voorgestelde framework zal:
\begin{itemize}
\item Een gedetailleerd overzicht bieden van de vereiste technische competenties, met specifieke aandacht voor systeembeheer binnen z/OS, waaronder platform fundamentals, storagebeheer, batch processing en JCL, automation en scripting (REXX), systeemmonitoring, security, storage en workload management.

\item De relevante soft skills identificeren die noodzakelijk zijn om effectief te functioneren binnen een mainframe-omgeving, zoals probleemoplossend vermogen, analytisch denken, documentatievaardigheden en samenwerking binnen enterprise IT-teams.

\item Concrete, praktijkgerichte oefeningen en use cases bevatten die gericht zijn op realistische operationele scenario’s, zodat studenten praktische ervaring opdoen met TSO/ISPF commando's, REXX, SDSF en JCL.

\item Gepresenteerd worden via een proof-of-concept webapplicatie die het competentieframework op een interactieve en toegankelijke manier visualiseert, waarbij verschillende competentiedomeinen en niveaus (Level 1 en Level 2) duidelijk worden weergegeven voor zowel studenten als opleiders.
\end{itemize}
Het uiteindelijke doel is om de kloof tussen onderwijs en industrie te verkleinen door studenten beter voor te bereiden op een carrière als z/OS SysAdmin/SysProg en zo bij te dragen aan het aanpakken van het groeiende tekort aan mainframe-specialisten binnen de sector.
%Wat is het beoogde resultaat van je bachelorproef? Wat zijn de criteria voor succes? Beschrijf die zo concreet mogelijk. Gaat het bv.\ om een proof-of-concept, een prototype, een verslag met aanbevelingen, een vergelijkende studie, enz.

\section{\IfLanguageName{dutch}{Opzet van deze bachelorproef}{Structure of this bachelor thesis}}%
\label{sec:opzet-bachelorproef}

% Het is gebruikelijk aan het einde van de inleiding een overzicht te
% geven van de opbouw van de rest van de tekst. Deze sectie bevat al een aanzet
% die je kan aanvullen/aanpassen in functie van je eigen tekst.
Deze bachelorproef is als volgt opgebouwd:

In Hoofdstuk~\ref{ch:stand-van-zaken} wordt een analyse gegeven van de huidige situatie in de mainframe
industrie, bestaande competentiemodellen en opleidingsmethoden, hierbij wordt
ook verwezen naar de bestaande literatuur hierover.

In Hoofdstuk~\ref{ch:methodologie} wordt een gedetailleerde beschrijving gegeven van de gebruikte
onderzoeksmethoden en-technieken.

In Hoofdstuk~\ref{ch:proof-of-concept} wordt de ontwikkeling en implementatie van het competentie frameworkbeschreven, evenals de proof-of-concept webapplicatie.

In Hoofdstuk~\ref{ch:conclusie}, tenslotte, wordt de conclusie gegeven, een antwoord geformuleerd
op de onderzoeksvragen, en worden aanbevelingen gedaan voor toekomstig onderzoek en implementatie.

In Hoofdstuk~\ref{ch:interviews} wordt de ontwikkeling en implementatie van het competentie frameworkbeschreven, evenals de proof-of-concept webapplicatie.

Het succes van deze bachelorproef zal worden gemeten aan de hand van de volgende criteria:
\begin{itemize}
\item De volledigheid en actualiteit van het competentie framework
\item De praktische toepasbaarheid van de voorgestelde oefeningen
\item De gebruiksvriendelijkheid en functionaliteit van de webapplicatie
\item De mate waarin het framework aansluit bij de behoeften van zowel het onderwijs als de industrie
\end{itemize}
%De rest van deze bachelorproef is als volgt opgebouwd:

%In Hoofdstuk~\ref{ch:stand-van-zaken} wordt een overzicht gegeven van de stand van zaken binnen het onderzoeksdomein, op basis van een literatuurstudie.

%In Hoofdstuk~\ref{ch:methodologie} wordt de methodologie toegelicht en worden de gebruikte onderzoekstechnieken besproken om een antwoord te kunnen formuleren op de onderzoeksvragen.

% TODO: Vul hier aan voor je eigen hoofstukken, één of twee zinnen per hoofdstuk

%In Hoofdstuk~\ref{ch:conclusie}, tenslotte, wordt de conclusie gegeven en een antwoord geformuleerd op de onderzoeksvragen. Daarbij wordt ook een aanzet gegeven voor toekomstig onderzoek binnen dit domein.